\section{The frame inequality at the selected level prime}

Write $T_u=(D.\mathrm{mem}(u)).A.T$ for the prime already assigned
to the word $u$. Let $g_0,g_1$ be the two host primes, $v$ the
host root discriminant divided by two, $m$ the number of bottom
weights. The rank record requires even $m\ge2$ and defines
$r_S=m/2-1$, so $2r_S=m-2$ also for its natural-number arithmetic.
The scale in the frame comparison is
$\mathrm{kaS}(D)=\pi/(v g_0 g_1)$. Thus a lower bound on $T_u$
has to retain the dependence on \emph{both} host primes and on
$v$. The arithmetic window already contains these three factors.

% BEGIN TFR1
\paragraph{TFR1. The lower edge of the prime window with its host factors.}
\label{claim:TFR1}
Use the ranks $m,k,f,r_S$ of an original datum and set
\[
 A=\log(vg_0g_1/\pi),\qquad
 X_u=\frac{1}{fk}\log\frac{D_*}{d_*(u)},\qquad L=x^\epsilon.
\]
Use the hypotheses of QW6 on the original row and column at $u$,
including $k\le x^{C_*}$, and the positive prime-window data
$T_u\ge (D_*/d_*(u))^{1/(fk)}\Gamma_u^{-2/(fk)}$.
If $m^2\le k$, $L\ge1$ and all the $m$ middle weights at every
place satisfy $w_{j\rho}\ge e$, then
\[
 \log T_u\ \ge\ mA+m-2C_\epsilon L
                    -\frac{2\log\Gamma_u}{fk},
 \qquad C_\epsilon=30+\frac{2+2C_*}{\epsilon}.       \tag{TFR.1}
\]
In particular, if $m\ge 2C_\epsilon L+2$ and
$2\log\Gamma_u/(fk)\le2$, then
\[
 T_u\ge(vg_0g_1/\pi)^m,\qquad
 T_u\,\mathrm{kaS}(D)^{2r_S}\ge(vg_0g_1/\pi)^2\ge1.
                                                        \tag{TFR.2}
\]

\begin{proof}
The weights $w_{j\rho}$ are the interlacing entries called $e_j$
in the pre-level record and $\mathrm{wj}(D,j+1,\rho)$ in the
family. Keep the entire first term of (QW.17). For each pair
$(j,\rho)$ it contains $\log w_{j\rho}+\log(v/\pi)$.
There are exactly $f$ places and $m$ such entries per place.
The division by $f$ leaves $m\log(v/\pi)$ and at least $m$
from the middle weights. Adding $m(\log g_0+\log g_1)$ and
bounding $L+m^2L/k$ by $2L$ gives
$X_u\ge mA+m-2C_\epsilon L$. The exact lower endpoint of the
window subtracts $2\log\Gamma_u/(fk)$ on taking logarithms.
This proves (TFR.1). Under the additional bounds,
$\log T_u\ge mA$. Finally $2r_S=m-2$ and $A\ge0$ by $v\ge\pi$
and $g_i\ge1$. Exponentiating gives (TFR.2). Here $d_*(u)$ is
the stored size at the center of the window; the arithmetic column
dimension $D.\mathrm{duS}(u)$ does not enter this calculation.
\end{proof}
% END TFR1

% BEGIN TFR2
\paragraph{TFR2. Uniform frame and gap bounds along a supplied family.}
\label{claim:TFR2}
On every original family $\mathcal F$ with eventual
$\operatorname{vol}(Y_0)=\operatorname{massI}(D,D.R.r)$
at every index, (TFR.2) holds at every sufficiently late $x$,
for all indices $i\in\mathcal F.Ix(x)$ and words $u$ of
$D=D_{x,i}$. Independently of the volume identity, every supplied
family satisfies the original $\mathrm{GapHyp}(D_{x,i})$ at
every index beyond one family threshold.

\begin{proof}
Fix $0<\epsilon<1/52$ and let
$C_*=\max(5,\mathcal F.\mathrm{CYL}+1)$. DSZ4 verifies the
actual interlacing, endpoint and host bounds required by QW6 on
all words of the family. In particular
$4kx\Delta\le b_0<w_{0\rho}$ and $\Delta\ge1$.
The bottom coordinates increase and $b_j<w_{j\rho}$, so every
$w_{j\rho}\ge e$ for $x\ge1$ since $k\ge2$.
The rank window gives $m\ge x^{1/4}$, $L\ge1$ and $m^2L\le k$.
With $C_\epsilon$ as in (TFR.1), take
\[
 x\ge \max\{1,(2C_\epsilon+2)^{1/(1/4-\epsilon)}\}.
\]
Then $m\ge (2C_\epsilon+2)L\ge 2C_\epsilon L+2$.

For the mass terms, DSZ5 identifies the member's stored row
degrees with the exact Harish--Chandra polynomials in the
compact-dual measure. The supplied base-volume identity and the
family fields $\mathrm{FactDstar},\mathrm{FactDuStar}$ give
the two row and column expressions in (QW.14), with equality
before inserting the permitted row factor $(2/3)^f$.
The factorial, index and positive Euler product calculation of
QW1--QW6 therefore applies with these very $D_*,d_*(u)$.
The original $\mathrm{Gam\_eq}$ and rank-size bound, as in
DSZ9, give $0<\log\Gamma_u\le fk$. Take the maximum of the
displayed threshold, the volume threshold and those of DSZ4--DSZ5
and DSZ9. TFR1 now proves the assertion for the prime in the given
record at \emph{every} late index.

For the gap, write $n=D.R.r=k+m-1$ for the total row rank.
DSZ4 gives $k\ge3x$ eventually, so $n\ge1801$ eventually.
For every original place and adjacent bottom pair (DSZ.12) gives
\[
 b_{j+1}-b_j\ \ge\ x^{\epsilon/100} n^{19}f^5 L.
\]
Taking also $x\ge16^{100/\epsilon}$ ensures the right side is
at least $16n^{19}$. These are both clauses in the definition of
$\mathrm{GapHyp}$. The bounds used for this last assertion are
family rank and weight fields alone.
\end{proof}
% END TFR2

% BEGIN TFR3
\paragraph{TFR3. The closed frame event and the concrete selector.}
\label{claim:TFR3}
For any family class $\mathcal A$ with a supplied
$\mathrm{TSel}(\mathcal A)$, the original closed
$\mathrm{PX.TFrame}$ holds on $\mathcal A$, in particular at
\[
 (C_g,c_b,C_u,c_T,C_T,C_0)
   =(\mathcal F.\mathrm{CYL}+1,\tfrac14,3,\tfrac12,
                                    \mathcal F.\mathrm{CTu},3).
\]
Given also the compatible $\mathrm{SelS}$ of
$\mathrm{link\_alt}$, there is an original
$\mathrm{CSel}$ on $\mathcal A$, hence on $D_3(3,\mathcal A)$, consisting of concrete
packages at \emph{every} sufficiently late index. More generally a
supplied $\mathrm{SSelF}$ and the eventual base-volume identity
of TFR2 suffice for $\mathrm{CSelF}$ on a single family.

\begin{proof}
For a member of $\mathcal A$, let $(a,R)$ be its selected
$\mathrm{TSelF}$. For $x>\mathcal F.x_1$ and $x\ge a$ the field
$R_{x,i}.T.\theta.\mathrm{real}$ gives an original
$\mathrm{RealTied}$ for exactly $D_{x,i}$. RSL3 identifies its
$Y_0$ volume with the standard genus mass there. The threshold
$\max(a,\mathcal F.x_1+1)$ suffices for this input to TFR2.
The conclusion gives $\mathrm{TFrameAt}(D_{x,i})$ without an
admissibility premise. Thus it gives the implication required by
$\mathrm{PX.TFrame}$ for any six constants in its definition,
at the family's own $sC',sC_l$. An $x$ threshold supplied by
TFR2 works simultaneously for all $i$ and all words. In
$\mathrm{Closed}$ one may take $\epsilon_1=1$ and the zero
transported bound, since every family in the class already has
that eventual event, independently of its additional largeness test.
In particular restriction along $D_3(3,\mathcal A)\subseteq
\mathcal A$ gives precisely the frame premise of $\mathrm{link\_alt}$.

At any member of $\mathcal A$, the compatible $\mathrm{SelS}$
supplies spectral ties $(\mathrm{selS}_{x,i}).\mathrm{toSTies}$
starting at the \emph{same} $a$ as $R_{x,i}$. Take the maximum
of $a$, the rank-gap threshold and the frame threshold. Beyond
it TFR2 supplies both numerical inputs to SFS6 at every index.
Use its $\mathrm{ConcK}$ at the selected $X,Y$ for that very
index. The threshold and this dependent function of $x,i$ are
$\mathrm{CSelF}$. Applying the construction to every member gives
the asserted $\mathrm{CSel}$. For the single-family statement
use its supplied $\mathrm{SSelF}$ at its own threshold and the
given base-volume identity in the same argument. The
distinct-word estimate, diagonal terms and counts in the full
spectral comparison are additional inputs at their original scopes.
\end{proof}
% END TFR3

% BEGIN TFR4
\paragraph{TFR4. The frame row before a family is chosen.}
\label{claim:TFR4}
Fix the source parameters $h_L:\mathrm{LocalLanglandsGL}$ and
$p:\mathrm{Bl.Prm}(h_L)$. Then the original row assertion
\[
\mathrm{RowE}(\mathrm{Bl.IsGoodE}(p.s),\mathrm{Bl.IsLinked}(p),
 \mathrm{TFrameAt})
\]
holds. In fact the estimate uses only the
construction context and the fact that the point lies over its
pre-level data. This is the row quantified over all compatible
points and bundles \emph{after} one common $x$ threshold.

\begin{proof}
Fix $0<\epsilon<1/52$, a standard constant record $C$, and the
construction exponents $c_D,c_h\ge0$. Let $K$ be any context
with these values at scale $x$. It supplies $P=K.P$, the level
choice $K.\mathrm{Lv}$ and $K.\mathrm{ok}:\mathrm{ChOK}$ for its
original coordinates and words. In (TFR.1) use $C_*=5$, hence
$C_\epsilon=30+12/\epsilon$. QW7--QW9 verify the place, rank
and discriminant hypotheses uniformly before $K$ or a word is
given. In particular $k\le x^5$, $m^2\le k$, $v\ge\pi$,
$w_{j\rho}>b_0\ge4kx\Delta$ and $m\ge x^{1/4}$ eventually.

The two size equations in $K$ are precisely
\begin{align*}
 d_*(u)&=\mu_k I_{\rm col}\prod_\rho d_k(u_\rho),\\
 D_*&=\mu_{k+m-1} I_{\rm row}\prod_\rho
                  d_{k+m-1}(v_\rho,-b_\rho),
\end{align*}
where $d_t(z)=\prod_{i<j}(z_i-z_j)/\prod_{i=0}^{t-1}i!$.
For the second equation the source defines
$\mathrm{trueDsig}$ by the formal-degree \emph{polynomial}.
Its entries are exactly the appended $v,-b$ in decreasing order,
and collecting denominator factors by their larger index gives
$\prod_{j=0}^{k+m-2}j!$. The pre-level record's $v,b$ agree with
those of the base word. Thus the stated equation holds for every
word and implies the required row lower bound with $(2/3)^f$.
The two level indices here are exactly those counted in QW2,
with $f$ copies of each of the two split-prime factors. The
standard masses are the ones in MF19 and the telescope of QW3.
This verifies the arithmetic inputs to (TFR.1) on $K$ itself.

There is one difference in the window comparison on a raw
context. Its actual $\Gamma$ is
$\exp(C_\Gamma m^5 fL)$, and $K.\mathrm{Lv}.\mathrm{T\_win}$
applies to this value. The rank fields give
$m\le2\sqrt{x}$ because $\delta_1=x^{-\epsilon/8}\le1$,
and $k\ge x^{9/2+10\epsilon}\ge x^{9/2}$. Hence
\[
 \frac{2\log\Gamma}{fk}
       \le 64|C_\Gamma|\,x^{-2+\epsilon}.
\]
This is at most $2$ once
$x\ge\max(1,(32|C_\Gamma|)^{1/(2-\epsilon)})$.
Take also the threshold of TFR2's elementary comparison
$x^{1/4}\ge2C_\epsilon L+2$, and those of QW7--QW9.
TFR1 now gives (TFR.2) for each prime $K.\mathrm{Lv}.T(u)$.

For an actual point $\beta$ over $K$, its core field
$\beta.c.\mathrm{isOver}$ identifies the stage base with
$\mathrm{Pre.ofCh}(P,K.\mathrm{Lv},\ldots)$. This construction
takes its words and ranks from $P$ and its prime function from
$K.\mathrm{Lv}$. The enlarged level-prime set in
$K.\mathrm{Lv}.\mathrm{host}$ preserves $g_0,g_1$ and the
underlying field defining $v$. The stage's column field
$\mathrm{A\_T}(u)$ identifies the member's $T_u$ with the base
prime at $u$. Finally $\beta.D$ is obtained from this stage by
$\mathrm{setDiv}$, which updates the division lattices and
their heights. Projecting the words, $T_u$, ranks and host from
its definition gives the quantities just bounded. Thus (TFR.2)
is the demanded $\mathrm{TFrameAt}(\beta.D)$.

For the outer $\mathrm{RowE}$ quantifiers one can use
$\epsilon_0=1$, the zero transported bound, and lower thresholds
$c_1=c_2=0$ for $c_D,c_h$. After the standard $C$, $c_D$ and
$c_h$ have been given, choose the finite maximum of the thresholds
above. The comparison then holds on every $K$ with
$\kappa.\mathrm{Of}(K)$ at that scale, every eligible point,
each real $c$ and every linked bundle over it. This gives the
original row for the supplied $h_L,p$. On this construction
route the row uses the written polynomial $\mathrm{trueDsig}$.
On general families TFR2 instead uses DSZ5 to identify their
stored analytic degrees.
\end{proof}
% END TFR4
